\section*{Chapter \ref{ch:pl:pricing-library-architecture} --- Pricing-Library Architecture}

\begin{solution}{exo:pl:pricing-library-architecture:1}
\texttt{at expiry \{ if S > 105 \{ pay 10; \} \}} with the schedule \texttt{expiry} holding the expiry date. On the chapter's market it prices at
4.09, equal to the library's digital option priced by its Monte Carlo engine on the same paths.
\end{solution}

\begin{solution}{exo:pl:pricing-library-architecture:2}
The autocallable is compared with the library's Monte Carlo engine on the same paths (same identifier, same seed, same dates) and the same
payoff arithmetic: the numbers are identical. The European is compared with the analytic formula, which has no Monte Carlo error; the script's
estimate differs from it by its own sampling error, 0.82 standard errors here.
\end{solution}

\begin{solution}{exo:pl:pricing-library-architecture:3}
Because the term sheet pays the coupon on an observation date and then redeems if the note is called: with the autocall first, \texttt{stop}
would end the path before the coupon of the call date was paid, and the script would undervalue every called path by a coupon (and its memory).
\end{solution}

\begin{solution}{exo:pl:pricing-library-architecture:4}
It removes the paths under the current mask from the living set; every later statement and event runs under \texttt{mask \& alive}, so those
paths pay and change nothing more. Skipping the rest of the block would skip it for all paths, including those for which the \texttt{stop} did
not apply.
\end{solution}

\begin{solution}{exo:pl:pricing-library-architecture:5}
The compiled evaluator pays a fixed cost per event (walking the tree once) and a small cost per path; the interpreter pays the walk per path. At
few paths the fixed cost and path generation (common to both) weigh more; as paths grow the ratio approaches the ratio of per-path costs, and
levels off once both are dominated by array work and memory traffic.
\end{solution}

\begin{solution}{exo:pl:pricing-library-architecture:6}
Observers for the screen, where one market quote changes at a time and only the prices on screen need recomputing; snapshots for the risk run,
where thousands of bumped markets are evaluated independently and in parallel, and nothing should depend on the order of updates.
\end{solution}

\begin{solution}{exo:pl:pricing-library-architecture:7}
97.64 at 50\%, 94.25 at 60\%, 92.28 at 70\%: about 5.4 points over twenty points of barrier. A lookback knock-in is very sensitive to the
barrier, since any observation below it counts; the barrier level is a pricing parameter to negotiate and a risk to hedge.
\end{solution}

\begin{solution}{exo:pl:pricing-library-architecture:8}
Independent random numbers make the script and the library differ by Monte Carlo noise, which hides small payoff errors; and they give up
common random numbers, so bumped script prices are noisy and Greeks unstable. Draw the library's paths with the instrument's seed; cross-check
by comparing with engines on the same paths (exact) and with closed forms (within the standard error).
\end{solution}

\begin{solution}{pb:pl:pricing-library-architecture:1}
\begin{enumerate}
\item Instruments, \omterm{def:pl:pricing-library-architecture:market}{market objects}, models and engines; the repricing function on snapshots connects them.
\item Questions about one piece of market state (a discount factor, a volatility); a snapshot never changes, so a bump makes a new one and
nothing is stale.
\item Common random numbers: the seed depends only on the instrument, so base and bumped prices use the same paths.
\item \texttt{supports} and \texttt{price} for an (instrument, model) pair.
\item Because every Greek, scenario and batch is that function applied to bumped snapshots.
\item State, events on schedules, assignments, conditions, payments and stops; \texttt{S}, \texttt{perf}, \texttt{t}, the parameters and the
state.
\item From the \texttt{at} blocks: every date of every block, sorted by date, in script order within a date.
\item The interpreter walks the tree per path with scalars; the compiled evaluator walks it per event with arrays and masks.
\item A state variable for the lowest performance, updated at every observation, and the maturity test on it.
\item The parser raises an error with its line and column.
\item European 11.3953 against 11.3485 (0.82); Asian 6.8012 against 6.8012 (0.00); down-and-out 11.0373 against 11.0226 (0.26);
autocallable 95.8750 against 95.8750 (0.00); lookback 94.2474, no library class.
\item Both are compared with the library's Monte Carlo engine on the same paths and payoffs.
\item 0.1518 per unit of spot per 100 of notional, both.
\item 18.8 times at 1\,000 paths, 33.4 at 30\,000, 34.3 at most (at 10\,000).
\item Two calculations for four states of the market: one at the first read, one at the read after three changes.
\item \textbf{Named result.} The scripted autocallable prices at 95.8750 per 100 with a delta of 0.1518, exactly as the library's autocallable
engine on common random numbers (0.00 standard errors apart); the compiled script runs 22 to 37 times faster than the interpreted one.
\item Validation against library engines and closed forms, review of the script, its storage with a version, and model approval.
\item When it needs a numerical method the script engine does not offer (early exercise by regression, a PDE), or when it trades in volume
and deserves a dedicated, faster engine.
\item In version control as text, with the trade referring to a script version and its parameters; every change reviewed and repriced.
\item New products as data, priced and risk-managed by the machinery the library already trusts.
\end{enumerate}
\end{solution}

\begin{solution}{iq:pl:pricing-library-architecture:1}
So that each varies independently: a new model prices old products, a new engine speeds up old models, a new product uses existing engines;
risk and scenarios work on any combination through one repricing function.
\iqlookfor{Separation of concerns, one repricing path.}
\end{solution}

\begin{solution}{iq:pl:pricing-library-architecture:2}
Write it as a \omterm{def:pl:pricing-library-architecture:script}{payoff script} on the library's scripting engine, validate against the closest products the library prices, check Greeks and
limits, and have it reviewed; a new engine can come later if volume justifies it.
\iqlookfor{Scripting and validation.}
\end{solution}

\begin{solution}{iq:pl:pricing-library-architecture:3}
Evaluate over all paths at once with arrays: conditions become masks, payments masked sums, termination a living set; generate paths once;
avoid per-path Python loops, and fall back to compiled code if the tree walk itself becomes the cost.
\iqlookfor{Vectorisation by masks.}
\end{solution}

\begin{solution}{iq:pl:pricing-library-architecture:4}
Results marked stale on input changes and recomputed when read; bad when reads are constant (no saving), in parallel batch risk (shared
mutable state), and when notification chains grow long and their order matters.
\iqlookfor{Observers against immutable snapshots.}
\end{solution}

\begin{solution}{iq:pl:pricing-library-architecture:5}
A small grammar (events on schedules, state, conditions, payments, stops), a parser with error positions, schedule extraction, an evaluator on
the library's paths with common random numbers, registration as an engine, a validation suite against library products, versioned storage and
review of scripts.
\iqlookfor{Small language, library paths, validation, versioning.}
\end{solution}
